% ============================================================
% Notas de pesquisa: avaliação de originalidade e esboço de proposta
% "Saturation-Routed MTL" — early exit por tarefa guiado por análise
% de convergência (documento de trabalho para re-síntese posterior)
% ------------------------------------------------------------
% Origem: discussão de ago/2026 sobre a continuação do artigo AAAI
% ("Information, Behavior, and Format: Where Syntax Lives in BERT
%  Depends on Who Is Asking" — layerwise_aaai_paper.tex).
% Ideia avaliada: supervisionar cada tarefa do MTL na camada onde sua
% informação satura (UPOS ~3, DEPREL ~8-9, HEAD topo), truncando o
% gradiente de cada tarefa na sua camada de ancoragem, e propor isso
% como paradigma de treinamento + early exit por tarefa na inferência.
%
% Sugestão de chaves .bib a criar (URLs nos comentários de cada item):
%   sogaard2016deep      — Søgaard & Goldberg, ACL 2016
%   hashimoto2017joint   — Hashimoto et al., EMNLP 2017 (JMT)
%   sanh2019hierarchical — Sanh et al., AAAI 2019 (HMTL)
%   xin2020deebert       — já existe em layerwise_aaai.bib
%   xin2021berxit        — Xin et al., EACL 2021 (BERxiT)
%   gao2023pfberxit      — Gao et al., Neurocomputing 2023 (PF-BERxiT)
%   li2021tokee          — Li et al., ACL 2021 (early exit token-level
%                          para sequence labeling)
%   he2021stemcell       — He & Choi, EMNLP 2021 (Stem Cell Hypothesis)
% ============================================================

% ─────────────────────────────────────────────────────────────────────────────
\section{Avaliação de Originalidade: Early Exit por Tarefa Guiado por Convergência}
\label{sec:proposta_saturation_routed}
% ─────────────────────────────────────────────────────────────────────────────

\subsection{A Ideia sob Avaliação}
\label{subsec:proposta_ideia}

Propor um paradigma de treinamento multi-tarefa em que \textbf{cada tarefa é supervisionada na
profundidade onde sua informação se concentra}, conforme medido pela análise por camada: o
cabeçote de UPOS seria acoplado à camada de saturação da tarefa (camada $\sim$3 no BERTimbau
Base), de modo que seu gradiente moldasse apenas as camadas inferiores e a supervisão
\emph{não fosse repassada} às camadas seguintes; DEPREL seria ancorado nas camadas 8--9; e
HEAD, cuja resolução se concentra no quarto superior da rede, permaneceria na camada final.
Na inferência, cada tarefa sairia na sua camada de ancoragem (\textit{early exit} por tarefa),
com a economia computacional correspondente.

\subsection{Veredito Sintético}
\label{subsec:proposta_veredito}

\textbf{A ideia bruta --- supervisão de tarefas de baixo nível em camadas inferiores --- não é
nova}: possui \textit{prior art} direto de 2016--2019, anterior ao BERT, e enquadrá-la como
``novo paradigma de treinamento'' seria vulnerável a rejeição imediata por revisores que
conheçam essa literatura. \textbf{A originalidade defensável está no critério de colocação, não
na arquitetura}: derivar a profundidade de ancoragem de cada tarefa \emph{empiricamente}, a
partir da análise de saturação (sondas + \textit{logit lens} + ablações causais), e demonstrar
que essa colocação \emph{depende do decoder e do regime de pré-treino} --- o que invalida a
colocação a priori usada por todo o \textit{prior art}. Originalidade estimada:
\textbf{incremental a moderada}; viável como continuação natural do artigo AAAI (o AAAI
diagnostica; este segundo artigo prescreve).

\subsection{\textit{Prior Art} Direto}
\label{subsec:proposta_priorart}

\begin{enumerate}
  \item \textbf{Søgaard \& Goldberg (2016)}~[sogaard2016deep]: supervisão de tarefas de baixo
        nível em camadas inferiores de BiLSTMs empilhadas --- POS na camada 1,
        \textit{chunking} na camada 3. É exatamente o mecanismo proposto (o gradiente do POS
        molda apenas as camadas inferiores), com colocação fixada a priori.
        % https://aclanthology.org/P16-2038/
  \item \textbf{Hashimoto et al. (2017), \textit{Joint Many-Task Model}}~[hashimoto2017joint]:
        POS $\rightarrow$ \textit{chunking} $\rightarrow$ \textit{parsing} de dependências em
        camadas sucessivamente mais profundas, com \textit{successive regularization} contra
        interferência catastrófica. Um ``paradigma de treinamento'' hierárquico por
        profundidade, proposto em 2017, pré-BERT.
        % https://arxiv.org/abs/1611.01587
  \item \textbf{Sanh et al. (2019), HMTL}~[sanh2019hierarchical]: a mesma hierarquia de
        profundidade aplicada a tarefas semânticas (NER, EMD, RE, correferência), publicada na
        própria AAAI.
        % https://arxiv.org/abs/1811.06031
  \item \textbf{Linha de \textit{early exit} para eficiência}: DeeBERT~[xin2020deebert],
        BERxiT~[xin2021berxit], PF-BERxiT~[gao2023pfberxit], PABEE, ConsistentEE; inclui
        \textit{early exit} em nível de \textit{token} para \textit{sequence labeling}
        (Li et al., ACL 2021)~[li2021tokee]. Campo saturado: ``early exit para UPOS''
        isoladamente não constitui contribuição.
        % survey/lista: https://github.com/falcon-xu/early-exit-papers
  \item \textbf{He \& Choi (2021), \textit{Stem Cell Hypothesis}}~[he2021stemcell]: em MTL com
        encoders Transformer, as tarefas competem pelas mesmas camadas e cabeças de atenção.
        É o trabalho contra o qual a proposta de ancoragem por tarefa precisa se posicionar:
        ancorar cada tarefa na sua camada é uma resposta possível ao dilema que eles apontam.
        % https://arxiv.org/abs/2109.06939
\end{enumerate}

\subsection{Contribuições Defensáveis (o que o \textit{Prior Art} Não Tem)}
\label{subsec:proposta_contribuicoes}

\begin{enumerate}
  \item \textbf{Colocação medida, não assumida (contribuição metodológica central).} Todo o
        \textit{prior art} fixa a profundidade de cada tarefa a priori, pela hierarquia
        linguística clássica. A proposta: um procedimento que mede a camada de saturação por
        configuração (sondas por camada + \textit{logit lens} + ablação causal) e ancora o
        cabeçote ali, com \textit{read-out} calibrado --- o experimento de cabeças congeladas
        do artigo AAAI já prova que a leitura calibrada na camada de saturação custa
        $\leq$0,4 ponto. O argumento contra a colocação a priori vem dos próprios resultados:
        a camada de convergência de UPOS \emph{se move com o decoder} (8,3 no linear vs.\ 6,1
        no biaffine) e com o pré-treino (mBERT adia todas as tarefas para as camadas 11--12).
        \emph{Não existe ``camada do POS'' universal} --- logo, qualquer colocação fixa está
        errada em geral.
  \item \textbf{A pergunta causal em aberto: o que a ancoragem faz com as tarefas superiores?}
        Se o gradiente de UPOS para na camada 3, as camadas 4--12 ficam livres da pressão de
        manter UPOS decodificável. Pela \textit{Stem Cell Hypothesis}, isso deveria reduzir
        interferência e beneficiar HEAD/DEPREL; pelo achado de \emph{formato} do artigo AAAI
        (o decoder remodela o encoder), a ancoragem pode também deslocar \emph{onde e como} a
        sintaxe se forma. Qualquer resultado --- melhora, piora ou neutralidade --- é um
        achado, pois conecta explicabilidade a decisão de projeto de treinamento com evidência
        causal. Nenhum dos trabalhos anteriores responde essa pergunta.
  \item \textbf{Um modelo, múltiplos pontos de operação (consequência de engenharia).} O mesmo
        \textit{parser} implantado serve requisições apenas-UPOS saindo na camada de ancoragem
        (\textit{speedups} reais já medidos: 3,67$\times$ na camada 3; 2,32$\times$ na camada
        5, RTX~4090) e \textit{parsing} completo indo até o topo, sem manter dois modelos.
        Framing modesto isoladamente, forte como corolário dos itens 1--2.
\end{enumerate}

\subsection{Esboço de Desenho Experimental}
\label{subsec:proposta_experimentos}

\begin{itemize}
  \item \textbf{Claim central}: a colocação de cabeçotes em MTL deve ser derivada de análise
        de saturação por configuração, e não de hierarquia linguística a priori; estudo causal
        do efeito da ancoragem sobre as tarefas superiores.
  \item \textbf{Condições a comparar} (mesma grade $2 \times 2$ do artigo AAAI ---
        \{linear, biaffine\} $\times$ \{BERTimbau, mBERT\} --- que é o diferencial):
        \begin{enumerate}
          \item MTL padrão, todos os cabeçotes na camada final (modelos atuais --- baseline);
          \item colocação a priori à la Søgaard/JMT (POS baixo, parsing alto, camadas fixas);
          \item colocação pela saturação medida (a proposta);
          \item colocação aleatória/adversarial (controle --- separa o efeito da colocação
                certa do efeito de qualquer colocação);
          \item \textit{early exit} dinâmico por confiança (estilo BERxiT) como referência
                de eficiência.
        \end{enumerate}
  \item \textbf{Resultado-alvo}: mostrar que a colocação ótima \emph{muda} entre as células da
        grade (decoder $\times$ pré-treino) --- é o resultado que refuta a colocação a priori.
  \item \textbf{Métricas}: UPOS/DEPREL/UAS/LAS com bootstrap pareado (protocolo já
        estabelecido); perfil por camada pós-ancoragem (as sondas e o \textit{logit lens}
        re-aplicados aos modelos ancorados, para medir o deslocamento da formação da sintaxe);
        latência real por ponto de operação.
\end{itemize}

\subsection{Riscos e Mitigações}
\label{subsec:proposta_riscos}

\begin{itemize}
  \item \textbf{Resultado nulo}: se a ancoragem não alterar o desempenho das tarefas
        superiores, o artigo vira \textit{negative result} --- ainda publicável (e.g.,
        Findings), mas mais fraco. Mitigação: o claim metodológico (item 1) e o perfil por
        camada pós-ancoragem sustentam o artigo mesmo sob efeito nulo.
  \item \textbf{Uma língua}: revisores pedirão ao menos um \textit{treebank} UD adicional além
        do Porttinari para o claim de método.
  \item \textbf{Economia condicional}: em MTL, o \textit{early exit} por tarefa só economiza
        quando subconjuntos de tarefas são servidos isoladamente (HEAD exige a rede completa);
        ser explícito sobre isso, como já feito na Seção~\ref{sec:analise_camadas} da
        dissertação (analise\_camadas\_early\_exit.tex).
  \item \textbf{Posicionamento}: citar e diferenciar explicitamente
        [sogaard2016deep, hashimoto2017joint, sanh2019hierarchical, he2021stemcell] já na
        introdução; nunca apresentar a arquitetura como novidade --- \emph{vender o critério
        de colocação e o estudo causal, não a arquitetura}.
\end{itemize}

\subsection{Síntese para Re-avaliação}
\label{subsec:proposta_sintese}

O artigo AAAI diagnostica: \emph{onde a sintaxe mora depende de quem pergunta} (informação,
comportamento e formato dissociam; a profundidade de convergência move-se com decoder e
pré-treino). A proposta aqui avaliada prescreve: \emph{então meça a saturação antes de ancorar
cada tarefa, e eis o que acontece --- em desempenho, em custo e na organização interna do
encoder --- quando se ancora}. Originalidade da mecânica: baixa (2016--2019). Originalidade do
critério de colocação guiado por explicabilidade, da grade decoder$\times$pré-treino e da
pergunta causal sobre interferência: suficiente para um artigo completo, condicionada a
baselines fortes (colocação a priori e aleatória) e a pelo menos uma segunda língua.
